\section{What the efficacy evidence shows}\label{sec:evidence}
\subsection{The case for CBT}
Against placebo, CBT outperforms on anxiety target symptoms ($g=0.56$ across 41 trials; \citealt{carpenter2018cbt}).
Against inactive controls, effects are larger: a meta-analysis of 144 trials found $g=0.75$ (depression) to $0.88$ (social anxiety disorder), with $0.80$ for GAD and $0.81$ for panic disorder \citep{cuijpers2016update}.
A 2025 network of 375 trials gave anxiety-spectrum effects between $0.5$ and $1.0$, above $1.0$ for PTSD and specific phobia \citep{cuijpers2025unified}.
In social anxiety, individual CBT ranked best of the therapies compared in a network meta-analysis of 101 trials \citep{mayowilson2014sad}, and in GAD only CBT remained significantly better than treatment as usual at 3--12 months in a network of 65 trials \citep{papola2024gad}.
Any critique has to explain why CBT works at all.

\subsection{The headline shrinks with better controls and bias adjustment}
\begin{itemize}
\item \textbf{Control type.} More than 80\% of anxiety trials used waitlist controls, and the few that did not showed much smaller effects \citep{cuijpers2016update}; only 17.4\% of trials in that analysis were high quality.
In the 375-trial analysis, waitlist-controlled effects exceeded $g=0.94$, while care-as-usual comparisons gave $0.22$--$1.13$ \citep{cuijpers2025unified}.
\item \textbf{Bias.} In depression, psychotherapy effects fell from $g=0.70$ to $g=0.31$ after adjusting for several sources of bias, with only 23\% of studies at low risk of bias \citep{cuijpers2019eysenck}.
In a cohort of National Institutes of Health-funded depression trials, 13 of 55 initiated grants did not lead to publication \citep{driessen2015publicationbias}.
Adding six unpublished studies to 20 published control comparisons changed the pooled effect from $g=0.52$ (95\% CI $[0.37,0.68]$) to $0.39$ ($[0.08,0.70]$).
The reported 25\% decrease is a change in point estimates; the published--unpublished contrast was uncertain ($p=.07$), and the pooled treatment effect remained positive.
An umbrella review of 102 recent meta-analyses (3{,}782 RCTs) found most target-symptom effects small and pooled psychotherapy effects of $0.34$ against placebo or treatment as usual \citep{leichsenring2022umbrella}.
\item \textbf{Response and remission.} About half of anxiety patients respond or remit \citep{loerinc2015response,springer2018remission}.
\end{itemize}

\subsection{Evidence volume, allegiance and reviewer selection}\label{sec:selection}
Evidence volume, comparative efficacy and component causality are different questions.
A large literature can increase confidence in a treatment, but its size alone establishes neither superiority to active alternatives nor which component causes improvement.
Researcher allegiance---investigators' preference for a treatment---was associated with comparative outcomes across 30 meta-analyses \citep{munder2013allegiance}.
That association warrants scrutiny, but does not alone separate bias from differences in treatment efficacy, expertise or implementation.
Nor does allegiance adjustment necessarily eliminate a CBT advantage: one review retained an advantage after adjustment, with superiority over psychodynamic but not interpersonal or supportive therapies \citep{tolin2010cbt}.
These findings justify examining particular comparisons, not declaring all therapies equivalent.

Publication bias must also be distinguished from reviewer preference.
Among the 13 unpublished studies identified by \citet{driessen2015publicationbias}, two were known to have been submitted and rejected, six were never submitted, and submission status was unclear for five.
This is not evidence that CBT-oriented reviewers excluded competing therapies.
A randomized experiment did find that peer reviewers favoured a positive-result manuscript over an otherwise comparable no-difference version, but it involved orthopaedic journals, not psychotherapy \citep{emerson2010peerreview}.
We did not identify direct evidence establishing CBT-specific reviewer suppression in the sources examined for this extension.
Testing that hypothesis requires reviewer-orientation data and comparable manuscripts, with non-submission, editorial decisions and review recommendations distinguished.
Prospective registration, complete reporting, credible active comparators and independent replication address evidence selection without presuming that one therapeutic orientation is uniquely biased.

\subsection{Does the effect shrink over time?}
This is contested, and the evidence is mixed.
In depression, one meta-analysis reported a steady decline in CBT effects since its introduction \citep{johnsen2015decline}, but a reanalysis restricted to randomised comparisons judged the decline ``most likely a spurious finding'' \citep{cristea2017notfalling}, and a 409-trial meta-analysis found no reduction by publication year \citep{cuijpers2023cbtdepression}.
In anxiety, a meta-regression of 49 controlled studies over 30 years found a flat trend (publication year $B=-0.008$, $p=.24$): effects ``did not increase over the last 3 decades'' \citep{hofmann2025effectsizes}.
Placebo-controlled anxiety estimates have been reported at lower values in successive reviews (0.73, 0.56, 0.24), but those come from separate reviews with different case mixes, not from one meta-regression.
We therefore do \emph{not} claim a decline over time; the defensible statement is that effects have not grown and that they are sensitive to control condition and bias.
